\chapter{Análisis del comportamiento digital infantil}
\label{chap:comportamiento-digital}

\insertminitoc
\parindent1.5em

\section{Conceptualización y evolución del comportamiento digital infantil}

El concepto de "comportamiento digital" en la infancia ha trascendido la simple medición del tiempo de pantalla para abarcar un espectro multidimensional de interacciones, consumos y creaciones en el ciberespacio. Históricamente, el análisis de la conducta tecnológica infantil se centraba en la recepción pasiva de medios unidireccionales, como la televisión. Sin embargo, en el ecosistema hiperconectado actual, el comportamiento digital se define como el conjunto de acciones cognitivas, sociales y emocionales que el menor despliega al interactuar con interfaces computacionales y redes globales \cite{holmarsdottir_integrative_2025}. Esta conceptualización moderna reconoce que el dispositivo móvil no es solo un canal de entretenimiento, sino una extensión de la identidad del menor, donde se construyen relaciones, se experimentan dinámicas de grupo y se desarrolla gran parte de la socialización primaria. Investigaciones de la OECD subrayan que comprender este comportamiento requiere un análisis ecológico que integre el contexto familiar, escolar y la arquitectura persuasiva de las propias plataformas \cite{noauthor_full_2025}.\\

La evolución de este comportamiento está intrínsecamente ligada a la reducción de la edad de inicio en el uso de la tecnología. Estudios pediátricos recientes alertan que los patrones de interacción digital comienzan a forjarse en la primera infancia, a menudo antes de que el menor desarrolle capacidades de lenguaje estructurado. Según \cite{bozzola_use_2022}, el uso temprano de medios digitales altera la forma en que los infantes procesan los estímulos visuales y auditivos, estableciendo una línea base de "necesidad de hiperestimulación" que condiciona su comportamiento en etapas posteriores. Cuando estos niños alcanzan la edad escolar, su comportamiento digital ya está profundamente arraigado en ciclos de retroalimentación inmediata, lo que dificulta la transición hacia tareas analógicas que requieren atención sostenida y paciencia.\\



El comportamiento digital también se caracteriza por una marcada plasticidad adaptativa; los menores modifican su conducta dependiendo de la plataforma que utilizan. Un niño puede mostrar un comportamiento colaborativo y de liderazgo en un entorno de videojuegos multijugador, y simultáneamente exhibir una conducta pasiva, ansiosa o de comparación social en aplicaciones de consumo de video vertical. Esta variabilidad dificulta la creación de perfiles conductuales estáticos. La investigación de \cite{tintori_childrens_2023} destaca que la "huella digital" del menor —el rastro de datos que deja a través de sus interacciones— es el reflejo más preciso de su comportamiento real, superando en fiabilidad a los cuestionarios de autoevaluación donde los adolescentes tienden a subestimar su nivel de dependencia o exposición al riesgo.\\

Finalmente, la conceptualización del comportamiento digital debe considerar el anonimato y la desinhibición como factores catalizadores de conductas atípicas. El "efecto de desinhibición en línea" explica por qué menores con un comportamiento dócil en el entorno físico pueden involucrarse en episodios de ciberacoso verbal o \textit{trolling} en el entorno virtual. Esta disonancia conductual ocurre porque la interfaz digital atenúa las señales sociales no verbales (como el llanto o las expresiones de dolor del interlocutor), reduciendo la respuesta empática inmediata \cite{zhu_cyberbullying_2021}. Por tanto, el estudio del comportamiento digital infantil exige herramientas de observación capaces de capturar no solo qué hace el niño, sino el contexto emocional y social en el que se desarrolla dicha acción.\\

\section{Patrones de uso y la formación de hábitos digitales}

La formación de hábitos digitales en la infancia no es un proceso accidental, sino el resultado de interacciones repetidas con plataformas diseñadas algorítmicamente para maximizar la retención. Un análisis longitudinal que comparó cohortes infantiles con una década de diferencia demostró un cambio sísmico en los patrones de uso: mientras que en el pasado predominaba el consumo mediático lineal y diversificado, el patrón actual está dominado por la inmersión interactiva en dispositivos móviles personales \cite{bohnert_emerging_2021}. Este cambio de paradigma ha trasladado el uso de la tecnología desde los espacios comunes del hogar (como el televisor en la sala) hacia la privacidad del dormitorio, lo que reduce drásticamente las oportunidades de mediación parental incidental y fomenta hábitos de consumo solitario y prolongado.\\

El patrón de uso más prevalente en la actualidad es el "consumo fragmentado continuo", caracterizado por sesiones de conexión muy breves pero extremadamente frecuentes a lo largo del día. Según estudios de \cite{serra_smartphone_2021}, los adolescentes pueden desbloquear sus teléfonos inteligentes decenas o incluso cientos de veces diarias para revisar notificaciones, impulsados por un mecanismo de recompensa de proporción variable similar al de las máquinas tragamonedas. Esta revisión compulsiva se cristaliza en un hábito inconsciente que interrumpe constantemente los procesos de consolidación de la memoria y el aprendizaje, generando una dependencia conductual donde la simple ausencia del dispositivo físico provoca cuadros de ansiedad aguda.\\



A nivel contextual, los patrones de uso varían significativamente según el entorno físico del menor. Investigaciones sobre el comportamiento en la escuela revelan que la disponibilidad de conectividad móvil ha desdibujado las fronteras entre el tiempo de ocio y el tiempo de estudio \cite{kopecky_behaviour_2021}. El uso subrepticio de dispositivos bajo el pupitre o durante los periodos de descanso escolar expone a los menores a dinámicas de riesgo sin la red de seguridad del entorno familiar. Para entender la magnitud de estos cambios conductuales, la siguiente tabla compara los hábitos de consumo mediático pasivo (tradicional) con los patrones de interacción digital interactiva, basándose en la literatura de \cite{bohnert_emerging_2021} y \cite{selak_problematic_2025}:

\begin{table}[h!]
\centering
\caption{Comparativa entre el consumo mediático tradicional y el hábito digital algorítmico}
\begin{tabular}{|p{3cm}|p{4.5cm}|p{4.5cm}|p{2cm}|}
\hline
\textbf{Dimensión} & \textbf{Consumo Tradicional (TV/PC Compartida)} & \textbf{Hábito Digital (Smartphone/Algoritmos)} & \textbf{Referencia} \\ \hline
Estructura de Tiempo & Sesiones largas y planificadas. & Fragmentado, ubicuo y constante. & \cite{serra_smartphone_2021} \\ \hline
Mecanismo de Selección & Búsqueda activa o programación fija. & Recomendación algorítmica pasiva (\textit{Auto-play}). & \cite{selak_problematic_2025} \\ \hline
Contexto Social & A menudo público y familiar. & Privado, solitario y nocturno. & \cite{bohnert_emerging_2021} \\ \hline
Feedback del Sistema & Unidireccional, sin recompensa social. & Bidireccional (Likes, notificaciones). & \cite{kopecky_behaviour_2021} \\ \hline
\end{tabular}
\end{table}

La consolidación de estos hábitos tiene repercusiones directas en la salud física del menor, siendo el desplazamiento del sueño el patrón colateral más documentado. El uso nocturno (\textit{vamping}) se ha normalizado como un hábito social de pertenencia, donde los menores sienten la obligación de mantenerse conectados para no ser excluidos de las conversaciones de su grupo de pares \cite{soriano-molina_association_2025}. Romper estos bucles de hábitos algorítmicos es extraordinariamente difícil mediante simples advertencias verbales; requiere intervenciones técnicas de diseño de fricción (\textit{friction design}), como bloqueos temporales o filtros de pantalla, que obliguen al cerebro del menor a pasar de un procesamiento automático a un procesamiento consciente sobre su propio comportamiento digital.\\

\section{Indicadores clínicos y conductuales del uso digital problemático}

La frontera entre un uso intensivo de la tecnología y un "Uso Problemático de Internet" (UPI) se define mediante la aparición de indicadores conductuales que interfieren negativamente con el desarrollo vital del menor. Aunque el UPI no está universalmente clasificado como una adicción en todos los manuales de diagnóstico psiquiátrico, la literatura coincide en que presenta marcadores clínicos paralelos a los trastornos de control de impulsos. \cite{selak_problematic_2025} define este estado como una incapacidad persistente para regular el consumo de medios digitales, acompañado de sentimientos de irritabilidad, disforia o agitación cuando el acceso a la red es restringido. Este síntoma de abstinencia digital es el primer gran indicador de que el comportamiento ha pasado de ser un hábito recreativo a una necesidad compulsiva que gobierna el estado de ánimo del infante.\\

Un indicador conductual clave del uso problemático es la "pérdida de control temporal" y la tolerancia. Los menores que desarrollan UPI requieren periodos cada vez más prolongados de inmersión en estímulos digitales hiper-estimulantes (como videos cortos de ritmo frenético) para alcanzar el mismo nivel de satisfacción o alivio del estrés. Esta escalada cuantitativa se correlaciona con un declive medible en las interacciones prosociales y un aislamiento físico de su red de apoyo analógica \cite{masri-zada_impact_2025}. Además, el comportamiento problemático suele ir acompañado de conductas de engaño; los menores desarrollan estrategias complejas para evadir la supervisión parental, como la instalación de aplicaciones bóveda (\textit{vault apps}) que ocultan contenido, o la creación de perfiles falsos (\textit{finstas}) para operar fuera del radar familiar \cite{dumaru_its_2024}.\\

El impacto psicológico de este comportamiento problemático se manifiesta de forma aguda a través de trastornos del sueño y cuadros de ansiedad. La investigación de \cite{soriano-molina_association_2025} demostró que los menores con uso problemático presentan niveles crónicos de $FoMO$ (miedo a perderse algo), manteniéndolos en un estado de alerta cognitiva permanente que imposibilita el descanso profundo. Esta privación del sueño retroalimenta el ciclo: un menor exhausto carece de los recursos ejecutivos necesarios para regular sus emociones, volviéndose más irascible frente a las figuras de autoridad y más susceptible a dinámicas de victimización, ciberacoso y comparación social destructiva en plataformas de imágenes \cite{martin-barrado_association_2025}.\\

Finalmente, los indicadores de uso problemático son a menudo síntomas de malestares psicológicos subyacentes. El entorno digital funciona frecuentemente como un mecanismo de afrontamiento disfuncional (\textit{coping mechanism}) para menores que sufren de soledad, TDAH, o problemas de integración escolar. Los depredadores en línea y los agresores cibernéticos identifican rápidamente a estos menores altamente activos y emocionalmente dependientes de la validación virtual \cite{page_psychological_2025}. Por consiguiente, los sistemas de protección infantil no deben limitarse a penalizar el uso excesivo, sino que deben estar capacitados para alertar a los padres cuando se detecte una agrupación de estos indicadores conductuales, facilitando una intervención temprana que aborde la raíz psicológica del problema en lugar de censurar únicamente la pantalla.\\

\section{Métricas de observación y evaluación técnica del comportamiento}

La medición objetiva del comportamiento digital infantil es el núcleo tecnológico de cualquier sistema de supervisión inteligente. Los métodos tradicionales basados en autoinformes o encuestas han demostrado ser inexactos, ya que los menores tienden a subestimar drásticamente su tiempo real de pantalla y rara vez reportan incidentes de exposición a riesgos por temor a la confiscación del dispositivo \cite{alsoubai_profiling_2024}. Para lograr una supervisión efectiva, es imperativo transitar hacia métricas de observación técnica pasiva e ininterrumpida. Según el análisis de \cite{alqahtani_childrens_2023}, las aplicaciones modernas de seguridad deben registrar métricas cuantitativas granulares, como la frecuencia de desbloqueo del dispositivo, la duración de las sesiones por aplicación (\textit{App Usage Stats \gls{api}}) y los volúmenes de tráfico de red en horarios nocturnos, proporcionando una radiografía temporal exacta del hábito digital.\\

Sin embargo, las métricas puramente temporales son insuficientes para evaluar la peligrosidad del comportamiento. Pasar tres horas programando o leyendo un libro electrónico no tiene el mismo impacto neuropsicológico que pasar tres horas deslizando videos de forma pasiva o interactuando con desconocidos en foros anónimos. Por ello, la investigación técnica avanza hacia el "perfilado del comportamiento" (\textit{Behavioral Profiling}), que utiliza algoritmos de aprendizaje automático para establecer una línea base de la actividad "normal" del menor y generar alertas cuando se detectan desviaciones significativas \cite{alsoubai_profiling_2024}. Para comprender esta evolución metodológica, la siguiente tabla contrasta las métricas tradicionales de supervisión con las métricas avanzadas impulsadas por Inteligencia Artificial:

\begin{table}[h!]
\centering
\caption{Evolución metodológica en las métricas de observación del comportamiento digital}
\begin{tabular}{|p{3.5cm}|p{4cm}|p{4cm}|p{2cm}|}
\hline
\textbf{Categoría de Métrica} & \textbf{Supervisión Tradicional} & \textbf{Supervisión Avanzada (\gls{ia} / Edge)} & \textbf{Ref. Técnica} \\ \hline
Tiempo y Frecuencia & Cronómetro global de uso de pantalla. & Fragmentación de sesiones y horarios de \textit{vamping}. & \cite{alqahtani_childrens_2023} \\ \hline
Evaluación de Contenido & Lista negra de URLs / Filtros DNS. & Análisis visual de pantalla y Detección de Objetos. & \cite{rojas_technical_2025} \\ \hline
Análisis de Interacción & Historial de aplicaciones abiertas. & \gls{nlp} para análisis de sentimiento e intención en texto. & \cite{razi_sliding_2023} \\ \hline
Alerta de Riesgo & Basado en reglas rígidas (Ej: Bloqueo a las 22:00). & Alertas contextuales basadas en anomalías conductuales. & \cite{alsoubai_profiling_2024} \\ \hline
\end{tabular}
\end{table}

La captura de contexto es el avance más significativo en la observación del comportamiento. Herramientas basadas en Procesamiento de Lenguaje Natural (\gls{nlp}) permiten medir el "tono emocional" (\textit{Sentiment Analysis}) de las conversaciones entrantes y salientes, determinando si el menor está siendo sometido a un lenguaje agresivo o si está participando en dinámicas de acoso \cite{razi_sliding_2023}. Al combinar estas métricas de texto con el análisis continuo de fotogramas mediante visión artificial local (\textit{Edge AI}), el sistema puede correlacionar eventos; por ejemplo, si un pico en la frecuencia cardíaca (obtenido de wearables) coincide con la visualización de una imagen agresiva o con connotaciones sexuales en la pantalla del móvil. \\

El gran desafío de la recolección de estas métricas es la unificación de los datos. Los menores modernos viven en un ecosistema multidispositivo (teléfonos, tabletas, consolas de videojuegos), lo que fragmenta la visibilidad del comportamiento. Los estudios técnicos sugieren que las plataformas de supervisión más robustas son aquellas que logran centralizar estas métricas de observación sin agotar la batería del dispositivo primario \cite{koulaxidis_improving_2022}. Además, la presentación de estas métricas al tutor debe ser intuitiva; la "fatiga de datos" ocurre cuando los padres son bombardeados con logs técnicos indescifrables. La \gls{ia} debe interpretar estas métricas y presentarlas en la bitácora de evidencias como narrativas comprensibles y procesables.\\

\section{Implicaciones éticas en el perfilado y monitoreo de menores}

La capacidad tecnológica para observar, cuantificar y perfilar el comportamiento digital infantil abre un debate ético de gran envergadura. A medida que las aplicaciones de seguridad recolectan métricas cada vez más íntimas (como análisis de sentimiento en mensajes privados o capturas de pantalla), la línea entre la protección legítima y el estado de vigilancia invasivo se vuelve peligrosamente delgada. \cite{agbana_ethical_2025} advierten que perfilar algorítmicamente a los niños puede etiquetarlos injustamente o generar falsos positivos que resulten en castigos inmerecidos, lo que erosiona profundamente la confianza filial. La ética del diseño exige que los sistemas no solo sean precisos, sino que respeten la presunción de inocencia del menor, utilizando el perfilado conductual estrictamente como un mecanismo de alerta temprana para el padre y no como un reporte policial irrefutable.\\

La comercialización de estos perfiles conductuales es quizás el mayor riesgo sistémico actual. Investigaciones exhaustivas, como las realizadas por \cite{feal_angel_2020}, han revelado empíricamente que una multitud de aplicaciones de control parental gratuitas o comerciales exfiltran las métricas de comportamiento de los niños —incluyendo datos de localización, uso de apps y horarios de sueño— hacia \textit{brokers} de datos y redes publicitarias. Este rastreo clandestino expone a los menores a la construcción de una "huella digital persistente" (\textit{digital footprint}) que, según \cite{tintori_childrens_2023}, puede ser utilizada en el futuro por corporaciones para microsegmentación publicitaria hiper-personalizada, manipulando las vulnerabilidades psicológicas del niño con fines comerciales. \\



Para abordar estas fallas éticas, el diseño de herramientas de observación debe anclarse en el principio de "Privacidad por Diseño" (\textit{Privacy by Design}). Esto implica la adopción innegociable del procesamiento en el borde (\textit{Edge Computing}), garantizando que el análisis de métricas sensibles y la inferencia de la Inteligencia Artificial ocurran enteramente en el hardware del dispositivo móvil, sin que los datos crudos viajen a servidores remotos \cite{ali_getting_2023}. Si la arquitectura técnica asegura que el desarrollador de la aplicación no tiene acceso físico o criptográfico a la bitácora de evidencias, se elimina el riesgo de monetización de datos infantiles y se cumple rigurosamente con normativas internacionales de protección como la COPPA (Children's Online Privacy Protection Act).\\

Finalmente, la ética de la supervisión requiere la inclusión participativa del menor. Los estudios de co-diseño con adolescentes realizados por \cite{dumaru_its_2024} y \cite{mcnally_co-designing_2018} evidencian que los niños poseen un agudo sentido de la privacidad; toleran y hasta agradecen la supervisión cuando se les informa transparentemente sobre qué métricas se están recogiendo y bajo qué circunstancias se activan las alertas a sus padres. Un sistema ético es aquel que notifica al menor que está siendo protegido y le otorga la posibilidad de revisar su propia bitácora de eventos para aprender de su comportamiento digital. En este paradigma, la tecnología de monitoreo evoluciona de ser un instrumento de control coercitivo a convertirse en un andamiaje educativo que fomenta el desarrollo de una ciudadanía digital responsable.